% Canonical BandNet triatomic derivation.
% Requires generated/derivations/equations.tex.

This derivation extends the same free-standing, symmetric interaction model to
a three-mass unit cell. It groups interaction distances modulo three, constructs
the Bloch force matrix, and identifies the Hermitian dynamic matrix whose three
eigenvalues give the squared frequencies.

\subsection{Model and interaction groups}

Let positive masses $m_1$, $m_2$, and $m_3$ repeat on sites separated by $d$.
Cell $\ell$ contains displacements $u_\ell$, $v_\ell$, and $w_\ell$ at sites
$3\ell$, $3\ell+1$, and $3\ell+2$, and has period $L=3d$. Each site is connected
symmetrically to the sites $n$ positions away by a passive spring $k_n\geq0$,
$n=1,2,\ldots,N$, with no springs to ground. The site equation is
\begin{equation}
    \BandNetTriatomicEquationOfMotion.
    \label{eq:tri-site}
\end{equation}
Define
\begin{equation}
    \BandNetTriatomicPadding.
    \label{eq:tri-padding}
\end{equation}
This padding changes neither the forces nor the dispersion. For the $u_\ell$
site, group $j=1,2,\ldots,P$ has neighbors
\begin{equation*}
\begin{array}{c|cc}
\text{distance}&\text{right neighbor}&\text{left neighbor}\\\hline
3j&u_{\ell+j}&u_{\ell-j}\\
3j-2&v_{\ell+j-1}&w_{\ell-j}\\
3j-1&w_{\ell+j-1}&v_{\ell-j}
\end{array}
\end{equation*}
The relations for $v_\ell$ and $w_\ell$ follow by translating one or two sites.

\subsection{Bloch reduction and force matrix}

Use the cell-index Bloch gauge
\begin{equation}
\begin{aligned}
    u_\ell(t)&=Ue^{i(q\ell L-\omega t)},&
    v_\ell(t)&=Ve^{i(q\ell L-\omega t)},&
    w_\ell(t)&=We^{i(q\ell L-\omega t)},\\
    Q&=qL=3qd.
\end{aligned}
\label{eq:tri-bloch}
\end{equation}
For example, distance $3j-2$ contributes factors $e^{i(j-1)Q}$ toward $V$ and
$e^{-ijQ}$ toward $W$ in the $u_\ell$ equation. Distance $3j-1$ contributes
$e^{-ijQ}$ toward $V$ and $e^{i(j-1)Q}$ toward $W$. Collecting these factors and
the $-2k_n$ term from every spring gives
\begin{equation}
    \BandNetTriatomicGeneralizedEigenproblem.
    \label{eq:tri-generalized}
\end{equation}
The force matrix is
\begin{equation}
    \BandNetTriatomicForceMatrix,
    \label{eq:tri-force}
\end{equation}
with entries
\begin{equation}
    \BandNetTriatomicEntries.
    \label{eq:tri-entries}
\end{equation}
For real stiffnesses, $A$ is real and
\begin{equation}
    \BandNetTriatomicHermitianIdentities,
    \label{eq:tri-hermitian}
\end{equation}
so $\mathbf F(q)$ is Hermitian. The passive-spring energy representation implies
that the force operator is negative semidefinite. A generic formal proof of this
spectral claim is a recorded deferred obligation. The corresponding
mass-normalized dynamic matrix is
\begin{equation}
    \BandNetTriatomicDynamicMatrix.
    \label{eq:tri-dynamic}
\end{equation}
Its ordered eigenvalues are the three squared frequencies. This form preserves
Hermitian symmetry and is equivalent to the generalized eigenproblem for
positive masses.

\subsection{Five-interaction case}

For $N=5$, padding gives $P=2$, $N'=6$, and $k_6=0$. The entries become
\begin{equation}
    \BandNetTriatomicFiveEntries.
    \label{eq:tri-five}
\end{equation}
All five physical stiffnesses remain distinct; $k_6=0$ is bookkeeping only.

\subsection{Wave-number domain and consistency checks}

The first Brillouin zone is $|q|\leq\pi/(3d)$. The positive half is represented
by $\hat q=3qd/\pi\in[0,1]$.

\paragraph{Rigid translation.}
At $q=0$, let
$K_{\mathrm{cross}}=\sum_{j=1}^{P}(k_{3j-2}+k_{3j-1})$. Then
$A=-2K_{\mathrm{cross}}$ and $B=C=K_{\mathrm{cross}}$, so every row of
$\mathbf F(0)$ sums to zero. Thus $[1,1,1]^T$ lies in its kernel and produces a
zero squared frequency.

\paragraph{Nearest-neighbor limit.}
If only $k_1$ is nonzero, $A=-2k_1$, $B=k_1$, and $C=k_1e^{-iQ}$, giving the
standard nearest-neighbor three-site-cell force matrix.

\paragraph{Equal masses.}
If $m_1=m_2=m_3$, the three-site description is a zone-folded monoatomic chain
of spacing $d$.
